% Generated from the shared arithmetic article fragments.
\section{Why an absolute coprime pair-square estimate cannot close the short rows}
\label{sec:pair-barrier}
Keep the original sextic family
\[
 M_u(D)=D^{-1/2}\sum_n\mu(n)\nu(n)W(Nn/D)\chi_n(u),
 \qquad \chi_n(u)=(u/n)_6,
\]
with squarefree primary column ideals outside the fixed excluded set,
their full annular coefficients (fix $W$ supported in $[1,2]$), and all
nonzero Eisenstein element rows
$Nu\le H$. The cubic symbol is $\chi_n^2$, not $\chi_n$.
The following classical frame inequality \cite{WelchGeometry} gives a new
obstruction to one proposed estimate in this arithmetic program.

\begin{proposition}[Coprime pair-square obstruction]\label{prop:pair-barrier}
Let $A$ be an $R$-by-$N$ complex matrix with $|A_{ui}|=1$, and put
$K_{ij}=\sum_u A_{ui}\overline{A_{uj}}$. Then
\[
 \sum_{i\ne j}|K_{ij}|^2\ge NR(N-R).
\]
For the original character kernels, let $B_1(H,D)$ be the sum of
$|K_H(n,m)|^2$ over all distinct coprime squarefree column ideals in
$[D,2D]$. At $H=D^{3/4}$,
\[
 B_1(H,D)\gg \frac{D^{11/4}}{(\log D)^2}.
\]
In particular the proposed sufficient input
$B_1(H,D)\ll H^2D^\varepsilon$ fails for each fixed $\varepsilon<5/4$.
This does not refute the signed M\"obius mean square.
\end{proposition}
\begin{proof}
Set $F_{uv}=\sum_i A_{ui}\overline{A_{vi}}$. Reordering four finite sums
proves $\sum_{i,j}|K_{ij}|^2=\sum_{u,v}|F_{uv}|^2$.
The row diagonal contributes $RN^2$; the column diagonal contributes
$NR^2$. Removing the latter and retaining the former gives the inequality.

Restrict the column pairs in the nonnegative sum $B_1$ to distinct prime
ideals $p,q$ of norm in $[D,2D]$, outside the same fixed set. If $H<D$,
no such prime divides a nonzero row element: divisibility would imply
$Nu\ge Np$. Thus all original zero masks are exactly one on this ball,
$|\chi_p(u)|=1$, and the residual primitive finite character conductor is
$pq$, of norm between $D^2$ and $4D^2$. All these pairs have common divisor1.
The lattice count gives $R=(2\pi/\sqrt3)H+O(\sqrt H+1)$.
The classical unconditional fixed-field prime-ideal theorem gives
$N\sim D/\log D$; this is an external analytic input, not a GRH hypothesis.
At the stated endpoint $R=o(N)$, and the bound is therefore
$(1-o(1))N^2R\gg D^{11/4}/(\log D)^2$.
\end{proof}

To see the exact failed step, write $c_n=\mu(n)\nu(n)W(Nn/D)$ and
$C_D=\sum_n|c_n|^2=O(D)$. Cauchy over coprime pairs bounds the absolute
signed contribution by $D^{-1}C_DB_1^{1/2}$. The hoped-for pair-square
upper bound would have closed this particular route, but is false.
A large absolute pair-square mass is not a lower bound on the quadratic
form with the fixed $c_n$. Nor is the prime part a positive summand of the
full M\"obius energy: prime--composite cross terms can cancel. The original
low-gcd estimate and the positive arithmetic determinant bridge remain open.

\subsection{Weighted profiles, exact arithmetic examples and the surviving task}
The finite matrix proof also retains an arbitrary row profile. Suppose
$|A_{ui}|^2=w_u$ for every column, set $S=\sum_u w_u$ and
$Q=\sum_u w_u^2$. The row diagonals are $Nw_u$ and column diagonals are
$S$, so the same identity proves
\[
 \sum_{i\ne j}|K_{ij}|^2\ge N^2Q-NS^2.
\]
This is a finite statement. Applying it to an infinite Schwartz tail would
require its convergence and every newly encountered prime collision;
these are not discarded. If a weighted pair-square attempt includes the
annular $W$, one can instead choose a closed subannulus on which a nonzero
$|W|$ is bounded below. Its prime count still has order $D/\log D$.
The fixed $\nu$ has unit modulus outside its conductor. These facts do not
replace the original coefficients in a signed estimate.

Here are exact sextic computations with all nonzero elements of the
Eisenstein norm ball, and both prime ideals above each split rational prime
in the column annulus. These kernel diagnostics use $\nu=1$ and exclude
only primes above2,3; they do not compute a signed M\"obius energy.
Values are integer norms in $\Z[\omega]$.
\begin{center}\small
\begin{tabular}{rrrrrr}
$D$&$H$&$R$&$N$&$\sum_{i\ne j}|K_{ij}|^2$&$NR(N-R)$\\\hline
128&4&18&22&6840&1584\\
1024&9&36&138&651672&506736\\
4096&81&300&456&60080328&21340800\\
256&64&240&40&363672&$-1920000$
\end{tabular}
\end{center}
The final row has $H=D^{3/4}$ but $N<R$; its negative lower bound is
correctly vacuous. The asymptotic theorem follows from the proof, not an
extrapolation of these examples. For $p\equiv1\pmod3$ and
$w^2+w+1=0$ in $\mathbf F_p$, evaluate $a+b\omega$ as $a+bw$ and take
its $(p-1)/6$ power, lifting the six roots via $\zeta_6=1+\omega$.
Every zero is tested before a unit value is assigned. For instance
$1+3\omega$ has norm7 and vanishes at $\omega=2$ modulo7, so the
unimodularity hypothesis cannot be used beyond its stated domain.

\texttt{ShortRowPairBarrier.lean} verifies the finite frame identity,
row-diagonal bound, weighted inequality and unimodular specialization,
with actual complex sums and explicit hypotheses. It does not formalize
the ideal-character comparison or prime-ideal theorem. The deterministic
integer checker also verifies weighted profiles and256 scalar pair sums
independently of its matrix computation. No numerical model or learned
coefficient is used.

The next signed route must retain $\mu,\nu,W$ before Cauchy. An exact
starting identity in ideal Dirichlet convolution is
\[
 \mu=2a-a*a*\mathbf1+b*b*\mathbf1,\qquad
 a=\mu\,\mathbf1_{Nn\le U},\quad b=\mu-a.
\]
Indeed $\mu*\mathbf1=\delta$ implies
$(\mu-a)*(\mu-a)*\mathbf1=\mu-2a+a*a*\mathbf1$.
On squarefree $n$, each divisor factorization retains pairwise
coprimality; the original row zeros still multiply along it. Estimating
these type-I/type-II kernels with their true conductors is a new open
arithmetic task. The identity alone supplies no cancellation or GRH proof.
